% id: 24-21
% book: 베이직쎈 미적분1 · 01 함수의 극한
% section: 개념 12 함수의 극한의 도형에의 활용
% passage: 오른쪽 그림과 같이 원 $x^2+y^2=1$ 위의 점~$\pt{P}(t{,}\mkern3mu\mathopen{}\sqrt{1-t^2})$ $(0<t<1)$과 점~$\pt{A}(1{,}\mkern3mu\mathopen{}0)$을 지나는 직선이 $y$축과 만나는 점을 $\pt{B}$라 하자. 점~$\pt{P}$에서 $x$축에 내린 수선의 발을 $\pt{H}$라 할 때, 다음을 구하시오.
% figure: tikz:fig-24-21
% answer: 
% answer_source: 
% variant_level: 0 (원본 전사)
% 이 파일은 build.mjs 가 items.json 에서 만든다. 고칠 때는 items.json 을 고친다.
\smash{\raisebox{4.5mm}{\dmkichul{베이직쎈 미적분1 24쪽 21번}}}
두 점~$\pt{A}$, $\pt{P}$를 지나는 직선의 방정식\\ $\Rightarrow$ $y=\dfrac{\sqrt{1-t^2}}{\blank}(x-1)$
